\documentclass[11pt,a4paper]{article}
% Preâmbulo comum. Não duplicar pacotes nos módulos.
\usepackage[utf8]{inputenc}
\usepackage[T1]{fontenc}
\usepackage[brazil]{babel}
\usepackage{amsmath,amssymb,amsthm,mathtools}
\usepackage{geometry}
\geometry{a4paper,margin=2.3cm}
\usepackage{microtype}
\usepackage{enumitem}
\usepackage{graphicx}
\usepackage{xcolor}
\usepackage{booktabs}
\usepackage{hyperref}
\hypersetup{colorlinks=true,linkcolor=blue!40!black,urlcolor=blue!50!black}

\setcounter{tocdepth}{2}
\numberwithin{equation}{section}

\theoremstyle{plain}
\newtheorem{teorema}{Teorema}[section]
\newtheorem{proposicao}[teorema]{Proposição}
\newtheorem{lema}[teorema]{Lema}
\newtheorem{corolario}[teorema]{Corolário}

\theoremstyle{definition}
\newtheorem{definicao}[teorema]{Definição}
\newtheorem{exemplo}[teorema]{Exemplo}
\newtheorem{exercicio}{Exercício}[section]

\theoremstyle{remark}
\newtheorem{observacao}[teorema]{Observação}
\newtheorem{contraexemplo}[teorema]{Contraexemplo}
\newtheorem{intuicao}[teorema]{Intuição}

\newcommand{\R}{\mathbb{R}}
\newcommand{\N}{\mathbb{N}}
\newcommand{\Z}{\mathbb{Z}}
\newcommand{\Q}{\mathbb{Q}}
\DeclareRobustCommand{\eps}{\varepsilon}
\DeclareRobustCommand{\vet}[1]{\vec{#1}}
\DeclareRobustCommand{\dif}{\mathrm{d}}
\DeclareRobustCommand{\abs}[1]{\left\lvert #1 \right\rvert}
\DeclareRobustCommand{\norm}[1]{\left\lVert #1 \right\rVert}
\DeclareMathOperator{\grad}{grad}
\DeclareMathOperator{\Div}{div}
\DeclareMathOperator{\rot}{rot}

\newcommand{\base}[1]{\par\medskip\noindent\textbf{Base para relembrar.} #1\par\medskip}
\newcommand{\cuidado}[1]{\par\medskip\noindent\textbf{Quando a hipótese falha.} #1\par\medskip}

\newcommand{\figura}[2]{%
  \begin{figure}[htbp]
  \centering
  \IfFileExists{#1}{%
    \includegraphics[width=0.84\textwidth]{#1}%
  }{%
    \fbox{\parbox{0.8\textwidth}{\centering Figura ainda não gerada:\\ \texttt{#1}}}%
  }
  \caption{#2}
  \end{figure}%
}

\newenvironment{esqueleto}{%
  \par\medskip\noindent\textbf{Esqueleto da demonstração.}\par\nopagebreak
  \begin{enumerate}[leftmargin=1.4em]
}{\end{enumerate}\medskip}

\title{Técnicas de integração, integrais impróprias\\ e equações de primeira ordem}
\author{}
\date{}
\begin{document}
\maketitle
\tableofcontents

\base{Usamos o Teorema Fundamental do Cálculo, a regra do produto, a substituição \(u=g(x)\) com \(g\) de classe \(C^1\), e o fato de que uma função contínua num intervalo fechado e limitado é integrável. A integral indefinida denota antiderivadas: escrever \(\int f=F+C\) significa \(F'=f\). Admitimos de \(\R\) que toda função crescente e limitada superiormente possui limite finito.}

\section{Integração por partes}

A regra do produto diz que, se \(u\) e \(v\) são deriváveis,
\[
(uv)'=u'v+uv'.
\]
Se \(u\) e \(v\) são de classe \(C^1\) num intervalo, então \(u'\), \(v'\) e \((uv)'\) são contínuas, logo integráveis em qualquer subintervalo fechado e limitado.

\begin{teorema}[Integração por partes]
Sejam \(u\) e \(v\) de classe \(C^1\) em \([a,b]\). Então
\[
\int_a^b u(x)v'(x)\,\dif x
=\Bigl[u(x)v(x)\Bigr]_a^b-\int_a^b v(x)u'(x)\,\dif x.
\]
Na notação diferencial, com \(\dif u=u'\,\dif x\) e \(\dif v=v'\,\dif x\),
\[
\int u\,\dif v=uv-\int v\,\dif u,
\]
igualdade de antiderivadas num intervalo em que \(u\) e \(v\) são de classe \(C^1\).
\end{teorema}

\noindent\textbf{Demonstração.}
Integramos a regra do produto e aplicamos o Teorema Fundamental do Cálculo ao lado esquerdo:
\begin{align*}
\int_a^b (uv)'(x)\,\dif x
&=\int_a^b u'(x)v(x)\,\dif x+\int_a^b u(x)v'(x)\,\dif x,\\
\Bigl[u(x)v(x)\Bigr]_a^b
&=\int_a^b v\,\dif u+\int_a^b u\,\dif v.
\end{align*}
Isolando a segunda integral do segundo membro,
\[
\int_a^b u\,\dif v
=\Bigl[uv\Bigr]_a^b-\int_a^b v\,\dif u.
\]
A forma indefinida é a mesma identidade sem avaliar os extremos, acrescida da constante de integração.
\hfill\(\square\)

\begin{intuicao}
Partes transfere a derivada de um fator para o outro. Vale a pena quando \(v=\int\dif v\) é simples e \(\int v\,\dif u\) é mais fácil do que \(\int u\,\dif v\).
\end{intuicao}

\begin{exemplo}[A antiderivada de \(x e^{x}\)]
Seja \(u=x\) e \(\dif v=e^{x}\,\dif x\). Então \(\dif u=\dif x\) e \(v=e^{x}\). As duas funções são de classe \(C^1\) em \(\R\). Logo
\begin{align*}
\int x e^{x}\,\dif x
&=x e^{x}-\int e^{x}\,\dif x
=x e^{x}-e^{x}+C
=e^{x}(x-1)+C.
\end{align*}
Derivando para conferir: \(\bigl(e^{x}(x-1)\bigr)'=e^{x}(x-1)+e^{x}=x e^{x}\).
\end{exemplo}

\begin{exemplo}[A integral definida, com os extremos]
A mesma escolha, agora em \([0,1]\):
\begin{align*}
\int_0^1 x e^{x}\,\dif x
&=\Bigl[e^{x}(x-1)\Bigr]_0^1
=e^{1}(1-1)-e^{0}(0-1)
=0-1\cdot(-1)
=1.
\end{align*}
A constante da antiderivada desaparece na diferença \(F(1)-F(0)\). Os dois extremos entram na conta: em \(x=1\) o fator \((x-1)\) zera o termo, e em \(x=0\) sobra \(-e^{0}(0-1)=1\).
\end{exemplo}

\subsection{Fórmula de redução de \(\int\sin^{n}x\,\dif x\)}

\begin{teorema}
Seja \(n\ge 2\) um inteiro. Então, em \(\R\),
\begin{equation}
\int\sin^{n}x\,\dif x
=-\frac{\sin^{n-1}x\cos x}{n}+\frac{n-1}{n}\int\sin^{n-2}x\,\dif x.
\label{eq:reducao-seno}
\end{equation}
A igualdade é de antiderivadas. A potência \(n-2\) pode ser zero: \(\sin^{0}x=1\).
\end{teorema}

\noindent\textbf{Demonstração.}
A função \(\sin^{n-1}x\) é de classe \(C^1\) em \(\R\) porque \(n-1\ge 1\) é inteiro. Tomamos
\[
u=\sin^{n-1}x,\qquad \dif v=\sin x\,\dif x.
\]
Então
\[
\dif u=(n-1)\sin^{n-2}x\cos x\,\dif x,\qquad v=-\cos x,
\]
e
\begin{align*}
\int\sin^{n}x\,\dif x
&=-\sin^{n-1}x\cos x-\int(-\cos x)(n-1)\sin^{n-2}x\cos x\,\dif x\\
&=-\sin^{n-1}x\cos x+(n-1)\int\cos^{2}x\,\sin^{n-2}x\,\dif x.
\end{align*}
Usamos \(\cos^{2}x=1-\sin^{2}x\):
\begin{align*}
\int\cos^{2}x\,\sin^{n-2}x\,\dif x
&=\int(1-\sin^{2}x)\sin^{n-2}x\,\dif x
=\int\sin^{n-2}x\,\dif x-\int\sin^{n}x\,\dif x.
\end{align*}
Substituindo,
\[
\int\sin^{n}x\,\dif x
=-\sin^{n-1}x\cos x+(n-1)\int\sin^{n-2}x\,\dif x-(n-1)\int\sin^{n}x\,\dif x.
\]
Somamos \((n-1)\int\sin^{n}x\,\dif x\) aos dois lados:
\[
n\int\sin^{n}x\,\dif x
=-\sin^{n-1}x\cos x+(n-1)\int\sin^{n-2}x\,\dif x.
\]
Dividimos por \(n\), que é positivo, e obtemos \eqref{eq:reducao-seno}.
\hfill\(\square\)

\begin{exemplo}[A fórmula com \(n=2\)]
\[
\int\sin^{2}x\,\dif x
=-\frac{\sin x\cos x}{2}+\frac{1}{2}\int 1\,\dif x
=\frac{x}{2}-\frac{\sin x\cos x}{2}+C.
\]
Pela identidade \(\sin(2x)=2\sin x\cos x\), isto é \(\frac{x}{2}-\frac{1}{4}\sin(2x)+C\), que é a antiderivada obtida pelo ângulo duplo. A fórmula não foi apenas citada: o caso \(n=2\) reproduz uma conta independente.
\end{exemplo}

\begin{corolario}
Para \(n\ge 2\) inteiro,
\[
\int_0^{\pi/2}\sin^{n}x\,\dif x
=\frac{n-1}{n}\int_0^{\pi/2}\sin^{n-2}x\,\dif x.
\]
\end{corolario}

\noindent\textbf{Demonstração.}
O termo de fronteira da fórmula indefinida, avaliado em \([0,\pi/2]\), é
\[
\Bigl[-\frac{\sin^{n-1}x\cos x}{n}\Bigr]_0^{\pi/2}
=-\frac{\sin^{n-1}(\pi/2)\cos(\pi/2)}{n}+\frac{\sin^{n-1}(0)\cos(0)}{n}
=0,
\]
porque \(\cos(\pi/2)=0\) e \(\sin 0=0\). Sobra apenas o coeficiente \(\frac{n-1}{n}\).
\hfill\(\square\)

\section{Produtos de seno e cosseno}

Sejam \(m,n\) inteiros não negativos. Queremos \(\int\sin^{m}x\cos^{n}x\,\dif x\).

\begin{teorema}
\begin{enumerate}
\item Se \(m\) é ímpar, \(m=2k+1\) com \(k\ge 0\), então
\[
\sin^{m}x\cos^{n}x
=(1-\cos^{2}x)^{k}\cos^{n}x\cdot\sin x,
\]
e a substituição \(u=\cos x\), \(\dif u=-\sin x\,\dif x\), produz um polinômio em \(u\).
\item Se \(n\) é ímpar, separa-se um fator \(\cos x\) e usa-se \(u=\sin x\).
\item Se \(m\) e \(n\) são pares, as identidades
\[
\sin^{2}\theta=\frac{1-\cos(2\theta)}{2},\qquad
\cos^{2}\theta=\frac{1+\cos(2\theta)}{2},\qquad
\sin(2\theta)=2\sin\theta\cos\theta
\]
abaixam o grau até uma soma de cossenos de arcos múltiplos.
\end{enumerate}
Se \(m\) e \(n\) são ambos ímpares, os dois primeiros caminhos servem.
\end{teorema}

\begin{exemplo}[Um expoente ímpar]
\begin{align*}
\int\sin^{3}x\cos^{2}x\,\dif x
&=\int\sin^{2}x\cos^{2}x\sin x\,\dif x
=\int(1-\cos^{2}x)\cos^{2}x\sin x\,\dif x.
\end{align*}
Com \(u=\cos x\) e \(\dif u=-\sin x\,\dif x\),
\begin{align*}
\int\sin^{3}x\cos^{2}x\,\dif x
&=-\int(1-u^{2})u^{2}\,\dif u
=-\int(u^{2}-u^{4})\,\dif u
=-\frac{u^{3}}{3}+\frac{u^{5}}{5}+C\\
&=\frac{\cos^{5}x}{5}-\frac{\cos^{3}x}{3}+C.
\end{align*}
Conferência: a derivada é
\[
\cos^{4}x(-\sin x)-\cos^{2}x(-\sin x)
=\sin x\cos^{2}x(1-\cos^{2}x)
=\sin^{3}x\cos^{2}x.
\]
\end{exemplo}

\begin{exemplo}[Os dois expoentes pares]
\begin{align*}
\sin x\cos x&=\frac12\sin(2x),\\
\sin^{2}x\cos^{2}x
&=\frac14\sin^{2}(2x)
=\frac14\cdot\frac{1-\cos(4x)}{2}
=\frac{1-\cos(4x)}{8}.
\end{align*}
Portanto
\begin{align*}
\int\sin^{2}x\cos^{2}x\,\dif x
&=\int\frac{1-\cos(4x)}{8}\,\dif x
=\frac{x}{8}-\frac{\sin(4x)}{32}+C.
\end{align*}
A derivada \(\frac18-\frac{4\cos(4x)}{32}=\frac{1-\cos(4x)}{8}\) recupera o integrando.
\end{exemplo}

\section{Substituição trigonométrica}

Em todo este parágrafo, \(a>0\). O sinal da raiz é o ponto delicado: \(\sqrt{\,\cdot\,}\) denota a raiz não negativa, e \(\sqrt{w^{2}}=\abs{w}\).

\begin{teorema}[Três famílias]
Seja \(a>0\).
\begin{enumerate}
\item Para \(\sqrt{a^{2}-x^{2}}\) com \(x\in[-a,a]\), tome \(x=a\sin\theta\) e \(\theta\in[-\pi/2,\pi/2]\). Então \(\cos\theta\ge 0\), \(\dif x=a\cos\theta\,\dif\theta\) e
\[
\sqrt{a^{2}-x^{2}}=\sqrt{a^{2}\cos^{2}\theta}=a\cos\theta.
\]
No triângulo retângulo de hipotenusa \(a\), o cateto oposto a \(\theta\) é \(x\) e o cateto adjacente é \(\sqrt{a^{2}-x^{2}}\). Assim \(\sin\theta=x/a\) e \(\cos\theta=\sqrt{a^{2}-x^{2}}/a\).
\item Para \(\sqrt{a^{2}+x^{2}}\) com \(x\in\R\), tome \(x=a\tan\theta\) e \(\theta\in(-\pi/2,\pi/2)\). Então \(\sec\theta>0\), \(\dif x=a\sec^{2}\theta\,\dif\theta\) e
\[
\sqrt{a^{2}+x^{2}}=\sqrt{a^{2}\sec^{2}\theta}=a\sec\theta.
\]
No triângulo, o cateto adjacente a \(\theta\) é \(a\), o cateto oposto é \(x\) e a hipotenusa é \(\sqrt{a^{2}+x^{2}}\).
\item Para \(\sqrt{x^{2}-a^{2}}\) com \(x\ge a\), tome \(x=a\sec\theta\) e \(\theta\in[0,\pi/2)\). Então \(\tan\theta\ge 0\), \(\dif x=a\sec\theta\tan\theta\,\dif\theta\) e
\[
\sqrt{x^{2}-a^{2}}=\sqrt{a^{2}\tan^{2}\theta}=a\tan\theta.
\]
No triângulo, a hipotenusa é \(x\), o cateto adjacente a \(\theta\) é \(a\) e o cateto oposto é \(\sqrt{x^{2}-a^{2}}\).
\end{enumerate}
\end{teorema}

\cuidado{Fora do intervalo de \(\theta\) indicado, o cosseno ou a tangente pode ser negativo, e trocar \(\sqrt{w^{2}}\) por \(w\) sem o módulo introduz um sinal errado. No terceiro caso, \(x\le -a\) pede um intervalo em que \(\sec\theta\le -1\); tratamos isso no exemplo, derivando a fórmula candidata nos dois componentes \(\abs{x}>a\).}

\begin{exemplo}[Família \(\sqrt{a^{2}-x^{2}}\), com \(a=1\)]
Calculamos \(\int\sqrt{1-x^{2}}\,\dif x\) para \(x\in(-1,1)\). Com \(x=\sin\theta\), \(\theta\in(-\pi/2,\pi/2)\) e \(\dif x=\cos\theta\,\dif\theta\),
\[
\sqrt{1-x^{2}}=\cos\theta,
\]
\[
\int\sqrt{1-x^{2}}\,\dif x
=\int\cos^{2}\theta\,\dif\theta
=\int\frac{1+\cos(2\theta)}{2}\,\dif\theta
=\frac{\theta}{2}+\frac{\sin(2\theta)}{4}+C.
\]
Como \(\sin(2\theta)=2\sin\theta\cos\theta\),
\[
\frac{\theta}{2}+\frac{\sin\theta\cos\theta}{2}+C
=\frac12\arcsin x+\frac12 x\sqrt{1-x^{2}}+C,
\]
onde \(\cos\theta=\sqrt{1-x^{2}}\) porque \(\cos\theta\ge 0\) no intervalo escolhido. Nos extremos \(x=\pm 1\) o integrando se anula e a mesma expressão, contínua em \([-1,1]\), continua válida.

Derivada, para cravar o coeficiente:
\begin{align*}
\frac{\dif}{\dif x}\Bigl(x\sqrt{1-x^{2}}\Bigr)
&=\sqrt{1-x^{2}}+x\cdot\frac{-x}{\sqrt{1-x^{2}}}
=\frac{1-x^{2}-x^{2}}{\sqrt{1-x^{2}}}
=\frac{1-2x^{2}}{\sqrt{1-x^{2}}},\\[0.4em]
\frac{\dif}{\dif x}\Bigl(\frac12 x\sqrt{1-x^{2}}+\frac12\arcsin x\Bigr)
&=\frac12\cdot\frac{1-2x^{2}}{\sqrt{1-x^{2}}}+\frac12\cdot\frac{1}{\sqrt{1-x^{2}}}
=\frac{1-x^{2}}{\sqrt{1-x^{2}}}
=\sqrt{1-x^{2}}.
\end{align*}
\end{exemplo}

Antes do segundo exemplo, integramos \(\sec^{3}\theta\), que não é uma citação: sai de partes.

\begin{lema}
No intervalo \(\theta\in(-\pi/2,\pi/2)\), onde \(\sec\theta>0\) e \(\tan\theta\) tem o sinal de \(\theta\),
\[
\int\sec^{3}\theta\,\dif\theta
=\frac12\sec\theta\tan\theta+\frac12\ln\abs{\sec\theta+\tan\theta}+C.
\]
\end{lema}

\noindent\textbf{Demonstração.}
Sejam \(u=\sec\theta\) e \(\dif v=\sec^{2}\theta\,\dif\theta\). Então \(\dif u=\sec\theta\tan\theta\,\dif\theta\) e \(v=\tan\theta\), e
\begin{align*}
\int\sec^{3}\theta\,\dif\theta
&=\sec\theta\tan\theta-\int\tan\theta\cdot\sec\theta\tan\theta\,\dif\theta
=\sec\theta\tan\theta-\int\sec\theta\tan^{2}\theta\,\dif\theta.
\end{align*}
Como \(\tan^{2}\theta=\sec^{2}\theta-1\),
\begin{align*}
\int\sec^{3}\theta\,\dif\theta
&=\sec\theta\tan\theta-\int\sec\theta(\sec^{2}\theta-1)\,\dif\theta
=\sec\theta\tan\theta-\int\sec^{3}\theta\,\dif\theta+\int\sec\theta\,\dif\theta.
\end{align*}
A antiderivada de \(\sec\theta\) em \((-\pi/2,\pi/2)\) é \(\ln\abs{\sec\theta+\tan\theta}\): de fato,
\[
\frac{\dif}{\dif\theta}\ln\abs{\sec\theta+\tan\theta}
=\frac{\sec\theta\tan\theta+\sec^{2}\theta}{\sec\theta+\tan\theta}
=\frac{\sec\theta(\tan\theta+\sec\theta)}{\sec\theta+\tan\theta}
=\sec\theta.
\]
Logo
\[
2\int\sec^{3}\theta\,\dif\theta
=\sec\theta\tan\theta+\ln\abs{\sec\theta+\tan\theta}+C,
\]
e basta dividir por \(2\).
\hfill\(\square\)

\begin{exemplo}[Família \(\sqrt{a^{2}+x^{2}}\), com \(a=1\)]
Para \(x\in\R\), seja \(x=\tan\theta\) com \(\theta\in(-\pi/2,\pi/2)\). Então \(\dif x=\sec^{2}\theta\,\dif\theta\) e \(\sqrt{1+x^{2}}=\sec\theta\), de modo que
\[
\int\sqrt{1+x^{2}}\,\dif x
=\int\sec^{3}\theta\,\dif\theta
=\frac12\sec\theta\tan\theta+\frac12\ln\abs{\sec\theta+\tan\theta}+C.
\]
Voltando, \(\tan\theta=x\) e \(\sec\theta=\sqrt{1+x^{2}}\). Além disso \(\sqrt{1+x^{2}}+x>0\) para todo \(x\) real, porque \(\sqrt{1+x^{2}}>\abs{x}\ge -x\). O módulo pode sair:
\[
\int\sqrt{1+x^{2}}\,\dif x
=\frac12 x\sqrt{1+x^{2}}+\frac12\ln\bigl(x+\sqrt{1+x^{2}}\bigr)+C.
\]
\end{exemplo}

\begin{exemplo}[Família \(\sqrt{x^{2}-a^{2}}\), com \(a=1\), e o sinal]
Primeiro, \(x>1\). Seja \(x=\sec\theta\) com \(\theta\in(0,\pi/2)\). Então \(\tan\theta>0\), \(\sqrt{x^{2}-1}=\tan\theta\) e \(\dif x=\sec\theta\tan\theta\,\dif\theta\), logo
\begin{align*}
\int\sqrt{x^{2}-1}\,\dif x
&=\int\tan\theta\cdot\sec\theta\tan\theta\,\dif\theta
=\int\sec\theta\tan^{2}\theta\,\dif\theta
=\int\sec\theta(\sec^{2}\theta-1)\,\dif\theta\\
&=\int\sec^{3}\theta\,\dif\theta-\int\sec\theta\,\dif\theta
=\frac12\sec\theta\tan\theta+\frac12\ln\abs{\sec\theta+\tan\theta}-\ln\abs{\sec\theta+\tan\theta}+C\\
&=\frac12\sec\theta\tan\theta-\frac12\ln\abs{\sec\theta+\tan\theta}+C.
\end{align*}
Voltando para \(x\),
\[
\int\sqrt{x^{2}-1}\,\dif x
=\frac12 x\sqrt{x^{2}-1}-\frac12\ln\bigl(x+\sqrt{x^{2}-1}\bigr)+C,\qquad x>1.
\]
O argumento do logaritmo é positivo para \(x>1\).

Esta expressão também serve em \(x<-1\), mas isso não sai da substituição com \(\theta\in(0,\pi/2)\), intervalo em que \(\sec\theta\ge 1\). Conferimos por derivação em \(\{\abs{x}>1\}\), onde \(\sqrt{x^{2}-1}>0\):
\begin{align*}
\frac{\dif}{\dif x}\Bigl(x\sqrt{x^{2}-1}\Bigr)
&=\sqrt{x^{2}-1}+\frac{x^{2}}{\sqrt{x^{2}-1}}
=\frac{x^{2}-1+x^{2}}{\sqrt{x^{2}-1}}
=\frac{2x^{2}-1}{\sqrt{x^{2}-1}},\\[0.4em]
\frac{\dif}{\dif x}\ln\abs{x+\sqrt{x^{2}-1}}
&=\frac{1+\dfrac{x}{\sqrt{x^{2}-1}}}{x+\sqrt{x^{2}-1}}
=\frac{\sqrt{x^{2}-1}+x}{\sqrt{x^{2}-1}\,\bigl(x+\sqrt{x^{2}-1}\bigr)}
=\frac{1}{\sqrt{x^{2}-1}}.
\end{align*}
Na segunda linha, \(x+\sqrt{x^{2}-1}\neq 0\) para \(\abs{x}>1\): o produto \(\bigl(x+\sqrt{x^{2}-1}\bigr)\bigl(x-\sqrt{x^{2}-1}\bigr)=1\) impede que o primeiro fator se anule. Portanto
\begin{align*}
\frac{\dif}{\dif x}\Biggl(\frac12 x\sqrt{x^{2}-1}-\frac12\ln\abs{x+\sqrt{x^{2}-1}}\Biggr)
&=\frac12\cdot\frac{2x^{2}-1}{\sqrt{x^{2}-1}}-\frac12\cdot\frac{1}{\sqrt{x^{2}-1}}
=\frac{x^{2}-1}{\sqrt{x^{2}-1}}
=\sqrt{x^{2}-1}.
\end{align*}
A mesma fórmula é antiderivada em \((1,\infty)\) e em \((-\infty,-1)\). A constante pode ser uma em cada componente: os intervalos não se comunicam.
\end{exemplo}

\section{Frações parciais}

Seja \(f=P/Q\) uma função racional, com \(P\) e \(Q\) polinômios e \(Q\) não identicamente nulo. Se \(\deg P\ge\deg Q\), dividimos primeiro:
\[
\frac{P}{Q}=S+\frac{R}{Q},\qquad \deg R<\deg Q,
\]
com \(S\) e \(R\) polinômios. Só então decompomos \(R/Q\).

Sobre \(\R\), fatoramos \(Q\) em lineares e quadráticos de discriminante negativo. Para as três formas abaixo procuramos coeficientes reais, resolvemos o sistema e \emph{voltamos a multiplicar}: a identidade conferida prova a decomposição daquela função, sem apelo a um teorema geral não demonstrado.

\begin{itemize}
\item Fatores lineares distintos: \(\dfrac{A}{x-r}+\dfrac{B}{x-s}\).
\item Fator linear repetido, aqui de ordem \(2\): \(\dfrac{A}{x-r}+\dfrac{B}{(x-r)^{2}}\). Um fator simples ao lado entra com o seu próprio coeficiente.
\item Quadrático irredutível \(x^{2}+bx+c\), com discriminante \(b^{2}-4c<0\): \(\dfrac{Ax+B}{x^{2}+bx+c}\). A integral de \(\dfrac{x}{x^{2}+bx+c}\) produz logaritmo; o resto, depois de completar o quadrado, produz arco tangente.
\end{itemize}

\begin{exemplo}[Lineares distintos]
\begin{equation*}
\frac{3x+5}{(x-1)(x+2)}=\frac{A}{x-1}+\frac{B}{x+2}.
\end{equation*}
O grau do numerador é \(1\), menor que \(2\). Limpando o denominador,
\[
3x+5=A(x+2)+B(x-1).
\]
Em \(x=1\): \(8=3A\), logo \(A=\dfrac{8}{3}\). Em \(x=-2\): \(-1=-3B\), logo \(B=\dfrac{1}{3}\). Conferência:
\[
\frac83(x+2)+\frac13(x-1)=\frac{8x+16+x-1}{3}=\frac{9x+15}{3}=3x+5.
\]
Assim, para \(x\neq 1\) e \(x\neq -2\),
\begin{align*}
\int\frac{3x+5}{(x-1)(x+2)}\,\dif x
&=\frac83\int\frac{\dif x}{x-1}+\frac13\int\frac{\dif x}{x+2}
=\frac83\ln\abs{x-1}+\frac13\ln\abs{x+2}+C.
\end{align*}
\end{exemplo}

\begin{exemplo}[Linear repetido]
O grau de \(x+2\) é menor que o grau de \(x(x-1)^{2}\). Procuramos
\[
\frac{x+2}{x(x-1)^{2}}=\frac{A}{x}+\frac{B}{x-1}+\frac{C}{(x-1)^{2}}.
\]
Então
\[
x+2=A(x-1)^{2}+Bx(x-1)+Cx.
\]
Em \(x=0\): \(2=A\). Em \(x=1\): \(3=C\). Expandindo,
\[
A(x^{2}-2x+1)+B(x^{2}-x)+Cx=(A+B)x^{2}+(-2A-B+C)x+A.
\]
Comparamos com \(x+2\): o coeficiente de \(x^{2}\) dá \(A+B=0\), logo \(B=-2\). O coeficiente de \(x\) confirma: \(-2\cdot 2-(-2)+3=-4+2+3=1\), e o termo independente é \(A=2\). Portanto
\begin{align*}
\int\frac{x+2}{x(x-1)^{2}}\,\dif x
&=2\ln\abs{x}-2\ln\abs{x-1}+\int C(x-1)^{-2}\,\dif x\\
&=2\ln\abs{x}-2\ln\abs{x-1}-\frac{3}{x-1}+C_{\mathrm{int}},
\end{align*}
válida em cada intervalo que não contém \(0\) nem \(1\). A constante de integração foi escrita \(C_{\mathrm{int}}\) para não confundir com o coeficiente \(C\).
\end{exemplo}

\begin{exemplo}[Quadrático irredutível: logaritmo e arco tangente]
O polinômio \(x^{2}+1\) tem discriminante \(-4<0\). O numerador \(2x+1\) tem grau menor que \(3\). Procuramos
\[
\frac{2x+1}{(x^{2}+1)(x+1)}=\frac{Ax+B}{x^{2}+1}+\frac{C}{x+1},
\]
isto é,
\[
2x+1=(Ax+B)(x+1)+C(x^{2}+1).
\]
Em \(x=-1\): \(-1=C(1+1)\), logo \(C=-\dfrac12\). Expandindo,
\[
(A+C)x^{2}+(A+B)x+(B+C)=2x+1.
\]
Daí \(A+C=0\), \(A+B=2\) e \(B+C=1\). Com \(C=-\frac12\) sai \(A=\frac12\). Depois \(\frac12+B=2\), logo \(B=\frac32\). E \(\frac32-\frac12=1\), como exigido. Conferência rápida do numerador:
\[
\Bigl(\frac12 x+\frac32\Bigr)(x+1)-\frac12(x^{2}+1)
=\frac12 x^{2}+\frac12 x+\frac32 x+\frac32-\frac12 x^{2}-\frac12
=2x+1.
\]
Integramos, para \(x\neq -1\):
\begin{align*}
\int\frac{2x+1}{(x^{2}+1)(x+1)}\,\dif x
&=\int\frac{\frac12 x+\frac32}{x^{2}+1}\,\dif x-\frac12\int\frac{\dif x}{x+1}\\
&=\frac12\int\frac{x}{x^{2}+1}\,\dif x+\frac32\int\frac{\dif x}{x^{2}+1}-\frac12\ln\abs{x+1}.
\end{align*}
Na primeira integral, \(w=x^{2}+1\), \(\dif w=2x\,\dif x\), então \(\int\frac{x}{x^{2}+1}\,\dif x=\frac12\ln(x^{2}+1)\). Conclusão:
\[
\frac14\ln(x^{2}+1)+\frac32\arctan x-\frac12\ln\abs{x+1}+C.
\]
\end{exemplo}

\begin{exemplo}[O grau do numerador não é menor: dividir primeiro]
\[
\frac{x^{3}+x^{2}+x+2}{x^{2}+1}.
\]
A divisão polinomial:
\[
(x+1)(x^{2}+1)=x^{3}+x+x^{2}+1=x^{3}+x^{2}+x+1,
\]
\[
(x^{3}+x^{2}+x+2)-(x^{3}+x^{2}+x+1)=1.
\]
Logo
\[
\frac{x^{3}+x^{2}+x+2}{x^{2}+1}=x+1+\frac{1}{x^{2}+1},
\]
e
\[
\int\frac{x^{3}+x^{2}+x+2}{x^{2}+1}\,\dif x
=\frac{x^{2}}{2}+x+\arctan x+C.
\]
Sem a divisão, o numerador de uma ``fração parcial'' teria grau maior que o do quadrático, e o sistema usual não fecha.
\end{exemplo}

\section{Integrais impróprias}

\begin{definicao}[Tipo I]
Seja \(f\) contínua em \([a,\infty)\). Definimos
\[
\int_a^{\infty} f(x)\,\dif x
:=\lim_{b\to\infty}\int_a^{b} f(x)\,\dif x,
\]
quando esse limite existe e é finito. Nesse caso a integral \textbf{converge}. Se o limite é infinito ou não existe, a integral \textbf{diverge}. O caso \(\int_{-\infty}^{b}\) é análogo, com limite \(a\to-\infty\). A integral \(\int_{-\infty}^{\infty}\) converge quando, para algum \(c\in\R\), as duas integrais \(\int_{-\infty}^{c}\) e \(\int_{c}^{\infty}\) convergem; soma-se então os valores. Os dois limites são independentes.
\end{definicao}

\begin{definicao}[Tipo II]
Seja \(f\) contínua em \((a,b]\), possivelmente ilimitada quando \(x\to a^{+}\). Definimos
\[
\int_a^{b} f
:=\lim_{c\to a^{+}}\int_c^{b} f,
\]
se o limite é finito. Se a singularidade está em \(b\), usa-se \(\lim_{c\to b^{-}}\int_a^{c} f\). Se a singularidade está num ponto interior \(c\in(a,b)\), a integral imprópria converge somente quando \emph{os dois} limites
\[
\lim_{\eps\to 0^{+}}\int_a^{c-\eps} f
\qquad\text{e}\qquad
\lim_{\delta\to 0^{+}}\int_{c+\delta}^{b} f
\]
existem e são finitos, separadamente. O valor é a soma. Não é lícito forçar \(\eps=\delta\) para produzir um cancelamento.
\end{definicao}

\begin{teorema}[Potências no infinito]
Seja \(p\in\R\). A integral \(\int_1^{\infty} x^{-p}\,\dif x\) converge se, e somente se, \(p>1\). Nesse caso o valor é \(\dfrac{1}{p-1}\).
\end{teorema}

\noindent\textbf{Demonstração.}
Se \(p\neq 1\),
\[
\int_1^{b} x^{-p}\,\dif x
=\Bigl[\frac{x^{1-p}}{1-p}\Bigr]_1^{b}
=\frac{b^{1-p}-1}{1-p}
=\frac{1-b^{1-p}}{p-1}.
\]
Se \(p>1\), então \(1-p<0\), logo \(b^{1-p}\to 0\) quando \(b\to\infty\), e o limite é \(\dfrac{1}{p-1}\).

Se \(p<1\), então \(1-p>0\), logo \(b^{1-p}\to\infty\). O quociente \(\dfrac{b^{1-p}-1}{1-p}\) tem denominador positivo e numerador tendendo a \(+\infty\), portanto a integral diverge para \(+\infty\).

Se \(p=1\),
\[
\int_1^{b}\frac{1}{x}\,\dif x=\ln b-\ln 1=\ln b\to\infty.
\]
A integral diverge. Reunindo os três casos, há convergência exatamente quando \(p>1\).
\hfill\(\square\)

\begin{teorema}[Potências junto à origem]
Seja \(p\in\R\). A integral \(\int_0^{1} x^{-p}\,\dif x\) converge se, e somente se, \(p<1\). Nesse caso o valor é \(\dfrac{1}{1-p}\).
\end{teorema}

\noindent\textbf{Demonstração.}
Se \(p\le 0\), a função \(x^{-p}\) se estende continuamente a \([0,1]\): vale \(1\) se \(p=0\), e \(x^{-p}=x^{\abs{p}}\to 0\) se \(p<0\), de modo que o valor \(0\) em \(x=0\) a torna contínua. A integral de Riemann ordinária existe, e \(p\le 0\) implica \(p<1\). O cálculo do parágrafo seguinte também cobre este caso.

Se \(p\neq 1\), para \(0<a<1\),
\[
\int_a^{1} x^{-p}\,\dif x
=\Bigl[\frac{x^{1-p}}{1-p}\Bigr]_a^{1}
=\frac{1-a^{1-p}}{1-p}.
\]
Se \(p<1\), então \(1-p>0\) e \(a^{1-p}\to 0\) quando \(a\to 0^{+}\). O limite é \(\dfrac{1}{1-p}\).

Se \(p>1\), então \(1-p<0\) e \(a^{1-p}\to +\infty\). O numerador \(1-a^{1-p}\to-\infty\) e o denominador é negativo, logo o quociente tende a \(+\infty\). A integral diverge.

Se \(p=1\),
\[
\int_a^{1}\frac{1}{x}\,\dif x=-\ln a\to+\infty
\]
quando \(a\to 0^{+}\). Portanto há convergência exatamente quando \(p<1\).
\hfill\(\square\)

\begin{contraexemplo}[Valor principal não é integral imprópria]
A função \(1/x\) é contínua em \([-1,0)\) e em \((0,1]\). Calculemos os dois lados exigidos pela definição.
\begin{align*}
\int_{\eps}^{1}\frac{1}{x}\,\dif x
&=\ln 1-\ln\eps=-\ln\eps\to+\infty,\\
\int_{-1}^{-\eps}\frac{1}{x}\,\dif x
&=\ln\eps-\ln 1=\ln\eps\to-\infty,
\end{align*}
quando \(\eps\to 0^{+}\). Cada limite diverge, então \(\int_{-1}^{1}\frac{1}{x}\,\dif x\) diverge.

Se, à força, usarmos o mesmo \(\eps\) dos dois lados,
\[
\int_{-1}^{-\eps}\frac{1}{x}\,\dif x+\int_{\eps}^{1}\frac{1}{x}\,\dif x
=\ln\eps-\ln\eps=0.
\]
O limite desse simétrico existe e vale \(0\). Chama-se valor principal de Cauchy. Ele não está na definição: a definição pede dois limites finitos, e aqui nenhum dos dois é finito. Cancelar \(+\infty\) com \(-\infty\) não produz convergência.
\end{contraexemplo}

\section{Testes de comparação}

Os testes deste parágrafo exigem sinal. Sem isso, uma função oscilante pode ter integral convergente enquanto o módulo diverge.

\begin{teorema}[Comparação direta]
Sejam \(f\) e \(g\) contínuas em \([a,\infty)\), com \(0\le f(x)\le g(x)\) para todo \(x\ge a\).
\begin{enumerate}
\item Se \(\int_a^{\infty} g\) converge, então \(\int_a^{\infty} f\) converge e
\[
0\le\int_a^{\infty} f\le\int_a^{\infty} g.
\]
\item Se \(\int_a^{\infty} f\) diverge, então \(\int_a^{\infty} g\) diverge.
\end{enumerate}
\end{teorema}

\noindent\textbf{Demonstração.}
Sejam \(F(b)=\int_a^{b} f\) e \(G(b)=\int_a^{b} g\). Como \(f\ge 0\) e \(g\ge 0\), as funções \(F\) e \(G\) são crescentes, e \(0\le F(b)\le G(b)\).

Se \(\int g\) converge, \(G(b)\) tende a um limite finito \(L\) e, por ser crescente, \(G(b)\le L\). Então \(F(b)\le L\). Uma função crescente e limitada superiormente tem limite finito, digamos \(M\). Logo \(\int f\) converge e \(0\le M\le L\).

Se \(\int f\) diverge, o fato de \(F\) ser crescente impede que \(F\) tenda a um limite finito, logo \(F(b)\to+\infty\). Como \(G(b)\ge F(b)\), também \(G(b)\to+\infty\).
\hfill\(\square\)

\begin{teorema}[Comparação no limite]
Sejam \(f\) e \(g\) contínuas em \([a,\infty)\), com \(f(x)\ge 0\) e \(g(x)>0\) para todo \(x\ge a\). Suponha que
\[
\lim_{x\to\infty}\frac{f(x)}{g(x)}=L\in[0,\infty].
\]
\begin{enumerate}
\item Se \(0<L<\infty\), então \(\int_a^{\infty} f\) converge se, e somente se, \(\int_a^{\infty} g\) converge.
\item Se \(L=\infty\) e \(\int_a^{\infty} g\) diverge, então \(\int_a^{\infty} f\) diverge.
\end{enumerate}
O caso \(L=0\), com a hipótese de que \(\int g\) converge, fica para o exercício~A1. A demonstração é do mesmo tipo que a do item (i), com uma cota só de um lado.
\end{teorema}

\noindent\textbf{Demonstração.}
(i) Existe \(M\ge a\) tal que, para \(x\ge M\),
\[
\Bigl\lvert\frac{f(x)}{g(x)}-L\Bigr\rvert<\frac{L}{2}.
\]
Daí \(\frac{L}{2}<\frac{f(x)}{g(x)}<\frac{3L}{2}\), e, como \(g>0\),
\[
\frac{L}{2}\,g(x)<f(x)<\frac{3L}{2}\,g(x),\qquad x\ge M.
\]
As integrais de \(a\) até \(M\) são ordinárias, de funções contínuas, logo finitas. Pela comparação direta, a cauda \(\int_M^{\infty} f\) converge se a cauda de \(g\) converge, e a cauda de \(g\) converge se a cauda de \(f\) converge, porque \(g<\frac{2}{L}f\). Portanto as integrais a partir de \(a\) convergem ou divergem juntas.

(ii) Se \(f/g\to\infty\), existe \(M\ge a\) tal que \(f(x)/g(x)>1\) para \(x\ge M\), isto é, \(f(x)>g(x)\). Se a cauda de \(g\) diverge, a comparação direta empurra a divergência para \(f\).
\hfill\(\square\)

\cuidado{Sem uma desigualdade \(f\ge 0\), ou sem passar ao módulo, a comparação não vale. Uma primitiva pode oscilar e convergir, enquanto a integral do módulo diverge. O par clássico é \(\sin x/x\) e \(\abs{\sin x}/x\).}

\begin{teorema}[Convergência de \(\int_1^{\infty}\frac{\sin x}{x}\,\dif x\)]
A integral \(\int_1^{\infty}\frac{\sin x}{x}\,\dif x\) converge.
\end{teorema}

\noindent\textbf{Demonstração.}
Fixe \(b>1\) e integre por partes em \([1,b]\). As funções \(u(x)=1/x\) e \(v(x)=-\cos x\) são de classe \(C^1\). Com \(\dif v=\sin x\,\dif x\) e \(\dif u=-x^{-2}\,\dif x\),
\begin{align*}
\int_1^{b}\frac{\sin x}{x}\,\dif x
&=\Bigl[-\frac{\cos x}{x}\Bigr]_1^{b}-\int_1^{b}(-\cos x)\Bigl(-\frac{1}{x^{2}}\Bigr)\,\dif x\\
&=\Bigl(-\frac{\cos b}{b}+\cos 1\Bigr)-\int_1^{b}\frac{\cos x}{x^{2}}\,\dif x.
\end{align*}
O detalhe do sinal: \(\int v\,\dif u=\int_1^{b}(-\cos x)(-x^{-2})\,\dif x=\int_1^{b}\cos x\,x^{-2}\,\dif x\), e partes subtrai esta integral.

Quando \(b\to\infty\), \(\abs{\cos b/b}\le 1/b\to 0\), logo \(\cos b/b\to 0\).

Resta provar que \(\int_1^{\infty}\frac{\cos x}{x^{2}}\,\dif x\) converge. O módulo satisfaz
\[
0\le\frac{\abs{\cos x}}{x^{2}}\le\frac{1}{x^{2}},\qquad x\ge 1.
\]
A integral \(\int_1^{\infty} x^{-2}\,\dif x\) converge pelo teorema das potências, com \(p=2>1\), e vale \(1\). A comparação direta dá a convergência de \(\int_1^{\infty}\abs{\cos x}/x^{2}\,\dif x\). Em particular, \(H(b)=\int_1^{b}\cos x\,x^{-2}\,\dif x\) é de Cauchy no infinito: se \(b_{2}>b_{1}>1\),
\[
\abs{H(b_{2})-H(b_{1})}
\le\int_{b_{1}}^{b_{2}}\frac{\abs{\cos x}}{x^{2}}\,\dif x,
\]
e a cauda do módulo tende a zero porque a integral do módulo converge. Admitimos o critério de Cauchy para limites em \(+\infty\): uma função \(H\), definida em \([1,\infty)\), possui limite finito em \(+\infty\) se, e somente se, para todo \(\eps>0\) existe \(M\) tal que \(b_{2}>b_{1}\ge M\) implica \(\abs{H(b_{2})-H(b_{1})}<\eps\). Assim \(H(b)\) converge, e portanto
\[
\lim_{b\to\infty}\int_1^{b}\frac{\sin x}{x}\,\dif x
=\cos 1-\int_1^{\infty}\frac{\cos x}{x^{2}}\,\dif x
\]
existe e é finito.
\hfill\(\square\)

O exercício~A3 isola, sem este exemplo, o passo geral ``convergência absoluta implica convergência''.

\begin{teorema}[Divergência de \(\int_1^{\infty}\frac{\abs{\sin x}}{x}\,\dif x\)]
A integral \(\int_1^{\infty}\frac{\abs{\sin x}}{x}\,\dif x\) diverge.
\end{teorema}

\noindent\textbf{Demonstração.}
Para cada inteiro \(k\ge 1\), no intervalo \([k\pi+\pi/4,\,k\pi+3\pi/4]\) vale \(\abs{\sin x}\ge\sin(\pi/4)=\sqrt{2}/2\), e \(x\le (k+1)\pi\), logo \(1/x\ge 1/((k+1)\pi)\). O comprimento do intervalo é \(\pi/2\). Portanto
\begin{align*}
\int_{k\pi}^{(k+1)\pi}\frac{\abs{\sin x}}{x}\,\dif x
&\ge\int_{k\pi+\pi/4}^{k\pi+3\pi/4}\frac{\sqrt{2}/2}{(k+1)\pi}\,\dif x
=\frac{\sqrt{2}}{2}\cdot\frac{1}{(k+1)\pi}\cdot\frac{\pi}{2}
=\frac{\sqrt{2}}{4(k+1)}.
\end{align*}
Estes intervalos, para \(k\ge 1\), estão contidos em \([1,\infty)\) e são disjuntos salvo extremos. Somando de \(k=1\) até \(N\),
\[
\int_1^{(N+1)\pi}\frac{\abs{\sin x}}{x}\,\dif x
\ge\sum_{k=1}^{N}\frac{\sqrt{2}}{4(k+1)}.
\]
A série harmônica deslocada diverge, logo o segundo membro tende a \(+\infty\) quando \(N\to\infty\). A integral de \(1\) até \(b\) fica arbitrariamente grande e a integral imprópria diverge.
\hfill\(\square\)

\begin{observacao}
A integral de \(\sin x/x\) converge, e a do módulo diverge. Diz-se que \(\int_1^{\infty}\sin x/x\,\dif x\) converge condicionalmente. A primitiva oscila: partes produziu \(-\cos b/b\), que amortece, mais uma integral absolutamente convergente. O amortecimento usa o sinal alternado de \(\sin x\); o módulo destrói esse cancelamento e a comparação com a série harmônica aparece.
\end{observacao}

\section{Equações diferenciais de primeira ordem}

\subsection{Equação separável}

\begin{definicao}
Sejam \(g\) contínua num intervalo \(I\) e \(h\) contínua num intervalo \(J\). A equação
\[
y'=g(x)h(y)
\]
é separável. Uma solução em \(I\) é uma função \(y:I\to J\), de classe \(C^1\), que satisfaz a igualdade em todo ponto de \(I\).
\end{definicao}

\begin{teorema}
Nas hipóteses da definição:
\begin{enumerate}
\item Se \(h(y_{0})=0\), a função constante \(y(x)=y_{0}\) é solução em \(I\). Chama-se solução de equilíbrio.
\item Se \(y\) é solução e \(h(y(x))\neq 0\) para todo \(x\in I\), então
\[
\int\frac{1}{h(y)}\,\dif y=\int g(x)\,\dif x,
\]
no sentido de que uma antiderivada de \(1/h\), composta com \(y\), difere de uma antiderivada de \(g\) por uma constante.
\end{enumerate}
\end{teorema}

\noindent\textbf{Demonstração.}
(i) A derivada da constante é \(0\), e \(g(x)h(y_{0})=0\).

(ii) Pela regra da cadeia,
\[
\frac{\dif}{\dif x}\int^{y(x)}\frac{1}{h}
=\frac{y'(x)}{h(y(x))}
=\frac{g(x)h(y(x))}{h(y(x))}
=g(x),
\]
desde que \(h(y(x))\neq 0\). Duas antiderivadas de \(g\) diferem por uma constante.
\hfill\(\square\)

\cuidado{Dividir por \(h(y)\) apaga os equilíbrios, que são exatamente os zeros de \(h\). Depois de integrar, ou se reintroduz o equilíbrio à parte, ou se verifica se alguma constante da família geral o recupera. Pior: se \(h\) não é lipschitziana perto do zero, o equilíbrio pode não ser a única solução do problema de valor inicial. O exercício~A2 dá duas soluções explícitas, e uma terceira colada, para \(y'=3y^{2/3}\) com \(y(0)=0\).}

\begin{exemplo}[Equação, equilíbrio e condição inicial]
Resolvemos \(y'=2x(1-y)\) com \(y(0)=3\). Aqui \(g(x)=2x\), contínua em \(\R\), e \(h(y)=1-y\), contínua em \(\R\). O equilíbrio é \(h(y)=0\), isto é, \(y\equiv 1\). Ele não atende \(y(0)=3\).

Para uma solução com \(y\neq 1\),
\[
\int\frac{\dif y}{1-y}=\int 2x\,\dif x.
\]
A primeira antiderivada: com \(w=1-y\), \(\dif w=-\dif y\), sai \(-\ln\abs{1-y}\). Logo
\[
-\ln\abs{1-y}=x^{2}+C_{1},\qquad
\ln\abs{1-y}=-x^{2}+C,\qquad
\abs{1-y}=e^{C}e^{-x^{2}}.
\]
Escrevemos \(1-y=Ke^{-x^{2}}\) com \(K\neq 0\). O valor \(K=0\) devolve o equilíbrio \(y=1\), então a família
\[
y(x)=1-Ke^{-x^{2}},\qquad K\in\R,
\]
contém todas as soluções constantes e as que obtivemos ao separar. A condição \(y(0)=3\) dá \(3=1-K\), logo \(K=-2\), e
\[
y(x)=1+2e^{-x^{2}}.
\]
Conferência: \(y(0)=1+2=3\) e
\[
y'=-4xe^{-x^{2}},\qquad
2x(1-y)=2x\bigl(1-(1+2e^{-x^{2}})\bigr)=-4xe^{-x^{2}}.
\]
\end{exemplo}

\subsection{Equação linear}

\begin{definicao}
Sejam \(p\) e \(q\) contínuas num intervalo \(I\). A equação
\[
y'+p(x)y=q(x)
\]
é linear de primeira ordem. Uma solução é uma função \(y:I\to\R\) de classe \(C^1\) que satisfaz a igualdade em \(I\).
\end{definicao}

\begin{teorema}[Fator integrante]
Seja \(P\) uma antiderivada de \(p\) em \(I\), e seja \(\mu(x)=\exp\bigl(P(x)\bigr)\). Então \(\mu\) é de classe \(C^1\), \(\mu>0\), \(\mu'=p\mu\), e para toda \(y\) de classe \(C^1\)
\[
(\mu y)'=\mu\bigl(y'+p y\bigr).
\]
Em particular, \(y\) resolve a equação linear se, e somente se, \((\mu y)'=\mu q\), isto é,
\[
\mu(x)y(x)=\int\mu(x)q(x)\,\dif x.
\]
\end{teorema}

\noindent\textbf{Demonstração.}
A exponencial é positiva e de classe \(C^1\), e a regra da cadeia dá \(\mu'=e^{P}P'=p\mu\). Então
\[
(\mu y)'=\mu'y+\mu y'=p\mu\, y+\mu y'=\mu(y'+py).
\]
Se \(y'+py=q\), o segundo membro é \(\mu q\). Reciprocamente, se \((\mu y)'=\mu q\), dividimos por \(\mu>0\) e obtemos a equação. Integrar \((\mu y)'=\mu q\) é o Teorema Fundamental.
\hfill\(\square\)

\begin{exemplo}[Problema de valor inicial numérico]
Resolvemos \(y'+3y=6\) com \(y(0)=1\), em \(\R\). Aqui \(p(x)=3\) e \(q(x)=6\) são contínuas. O fator integrante é
\[
\mu(x)=\exp\Bigl(\int 3\,\dif x\Bigr)=e^{3x},
\]
tomando a constante da antiderivada de \(p\) igual a zero; outra constante só multiplicaria \(\mu\) por um positivo e cancelaria nos dois lados. A identidade do teorema dá
\[
\bigl(e^{3x} y\bigr)'=6e^{3x}.
\]
Integramos:
\[
e^{3x} y=6\cdot\frac{1}{3}e^{3x}+C=2e^{3x}+C,
\qquad
y=2+Ce^{-3x}.
\]
A condição inicial: \(1=2+C\), logo \(C=-1\), e
\[
y(x)=2-e^{-3x}.
\]
Conferência: \(y(0)=2-1=1\) e
\[
y'=3e^{-3x},\qquad
y'+3y=3e^{-3x}+3\bigl(2-e^{-3x}\bigr)=6.
\]
\end{exemplo}

\section{Exercícios}

\begin{exercicio}[Básico --- partes]
\textbf{B1.} Calcule \(\int x\cos x\,\dif x\) e confira por derivação.
\end{exercicio}

\begin{exercicio}[Básico --- expoente ímpar]
\textbf{B2.} Calcule \(\int\sin^{3}x\,\dif x\) separando um fator \(\sin x\).
\end{exercicio}

\begin{exercicio}[Básico --- substituição trigonométrica]
\textbf{B3.} Calcule \(\int\dfrac{1}{\sqrt{9-x^{2}}}\,\dif x\) para \(\abs{x}<3\), voltando à variável \(x\). Registre o sinal da raiz.
\end{exercicio}

\begin{exercicio}[Básico --- lineares distintos]
\textbf{B4.} Calcule \(\int\dfrac{1}{(x-2)(x+1)}\,\dif x\), resolvendo o sistema e conferindo a decomposição.
\end{exercicio}

\begin{exercicio}[Básico --- imprópria de tipo I]
\textbf{B5.} Calcule \(\int_1^{\infty} x^{-3}\,\dif x\) pelo limite. Diga o valor de \(p\) e por que a integral converge.
\end{exercicio}

\begin{exercicio}[Intermediário --- dois expoentes pares]
\textbf{I1.} Calcule \(\int\cos^{4}x\,\dif x\) pelas identidades de ângulo duplo, até uma combinação de \(x\), \(\sin(2x)\) e \(\sin(4x)\).
\end{exercicio}

\begin{exercicio}[Intermediário --- quadrático irredutível]
\textbf{I2.} Calcule \(\int\dfrac{x+3}{(x-1)(x^{2}+1)}\,\dif x\). O resultado deve conter logaritmo e arco tangente.
\end{exercicio}

\begin{exercicio}[Intermediário --- imprópria calculada]
\textbf{I3.} Calcule \(\int_0^{\infty}\dfrac{1}{1+x^{2}}\,\dif x\) pelo limite. A função é contínua em \([0,\infty)\); a única impropriedade está no infinito.
\end{exercicio}

\begin{exercicio}[Intermediário --- comparação no limite]
\textbf{I4.} Decida se \(\int_1^{\infty}\dfrac{\sqrt{x}}{x^{2}+1}\,\dif x\) converge. Use a comparação no limite com \(g(x)=x^{-3/2}\), escrevendo o valor de \(L\).
\end{exercicio}

\begin{exercicio}[Intermediário --- separável]
\textbf{I5.} Resolva \(y'=2x(y+1)\) com \(y(0)=0\). Trate o equilíbrio \(y\equiv -1\) antes de dividir, escreva a solução explícita do problema de valor inicial e confira por derivação.
\end{exercicio}

\begin{exercicio}[Intermediário --- linear com \(p\) não constante]
\textbf{I6.} Resolva \(y'+\dfrac{2}{x}y=x\) para \(x>0\), com \(y(1)=1\). Justifique o domínio \(x>0\) e confira a solução.
\end{exercicio}

\begin{exercicio}[Avançado --- comparação com \(L=0\)]
\textbf{A1.} Sejam \(f\) e \(g\) contínuas em \([a,\infty)\), com \(f\ge 0\), \(g>0\) e \(\lim_{x\to\infty} f(x)/g(x)=0\). Prove: se \(\int_a^{\infty} g\) converge, então \(\int_a^{\infty} f\) converge.
\end{exercicio}

\begin{exercicio}[Avançado --- equilíbrio sem unicidade]
\textbf{A2.} Considere \(y'=3y^{2/3}\) com \(y(0)=0\), a potência no sentido real \((y^{1/3})^{2}\).
\begin{enumerate}
\item Verifique que \(y(x)=0\) e \(y(x)=x^{3}\) são soluções em \(\R\).
\item Verifique que a função que vale \(0\) para \(x<0\) e \(x^{3}\) para \(x\ge 0\) também é solução.
\item Mostre que \(f(y)=3y^{2/3}\) não é lipschitziana em vizinhança alguma de \(y=0\), calculando o quociente \(\abs{f(y)-f(0)}/\abs{y}\).
\end{enumerate}
\end{exercicio}

\begin{exercicio}[Avançado --- absoluta implica convergente]
\textbf{A3.} Seja \(f\) contínua em \([a,\infty)\). Admita o critério de Cauchy: \(H:[a,\infty)\to\R\) tem limite finito em \(+\infty\) se, e somente se, para todo \(\eps>0\) existe \(M\) tal que \(b_{2}>b_{1}\ge M\) implica \(\abs{H(b_{2})-H(b_{1})}<\eps\). Prove que, se \(\int_a^{\infty}\abs{f}\) converge, então \(\int_a^{\infty} f\) converge.
\end{exercicio}

\end{document}
